\documentclass[luatex,fontsize=8pt,paper=b5,tate]{jlreq}
% \documentclass{ltjtarticle}
\usepackage{KKsymbols}
\usepackage{luwa-ul}
\usepackage{gckanbun,lltjext,luatexja-ruby}

\makeatletter

\begin{document}\fboxsep=0pt\fboxrule=.1pt%

\LARGE
\begin{GCKEnv}{2\zw}
  今夫\送り{レ}江戸\振り{者}{は}、世之所\送り{ノ}\返り{レ}称\送り{スル}名都\振り{大}{だい}\振り{邑}{いふ}、冠蓋之所\返り{レ}集\送り{マル}\LineNumbering*{\kakko{2}}\dashKKAuto{舟車之\振り{所}{ところ}\送り{ニシテ}\返り{レ}\振り{湊}{あつ}\送り[intrusion=post]{マル}}、実\送り{ニ}\振り{為}{た}\送り{ル}\返り{二}天下之大都会\返り{一}也。
  而\送り{レドモ}\LineNumbering{C}\underLineKKAuto{其地之為名、訪之於古、未之聞}。
  豈\送り{ニ}非\送り{ズ}\返り{三}古今相\送り{ヒ}去\送り{ルコト}\振り{日}{ひび}\送り{ニ}遠\送り{ク}、事之相\送り{ヒ}変\送り{ズルコト}愈多\送り{ク}、求\送り{ムルモ}\返り{二}其\送り{ノ}所\送り{ヲ}\返り{\IchiRe}欲\送り{スル}\返り{レ}聞\送り{カント}而不\送り{ルコト}可\送り{カラ}\返り{レ}得、亦\送り{タ}\振り{猶}{な}[ごと]\送り{ホ}[キヲ]\返り{二}今之於\送り{ケルガ}\返り{\JyouRe}古\送り{ニ}也。
\end{GCKEnv}

\bigskip

\begin{GCKEnv}{3\zw}
  \fbox{\GCKanshiBox{10\zw}{%
    \振り[intrusion=both]{春}{しゅん}%
    \振り[intrusion=both]{眠}{みん}%
    不\返り[intrusion=both]{レ}%
    \振り{覚}{おぼ}\送り{エ}\返り[intrusion=both]{レ}%
    \振り{暁}{あかつき}\送り[intrusion=post]{ヲ}%
  }}

  \GCKanshiBox{10\zw}{%
    \振り[intrusion=both]{処}{しょ}%
    \振り[intrusion=both]{処}{しょ}%
    聞\送り{ク}\返り[intrusion=both]{二}%
    \underLineKKAuto{\振り[intrusion=both]{啼}{てい}%
    鳥\送り{ヲ}\返り[intrusion=post]{一}}%
  }

  \GCKanshiBox{10\zw}{%
    \振り[intrusion=both]{夜}{や}%
    \振り[intrusion=both]{来}{らい}%
    風%
    雨\送り[intrusion=both]{ノ}%
    声%
  }

  \GCKanshiBox{10\zw}{%
    花%
    \振り{落}{お}\送り[intrusion=both]{ツルコト}%
    知\送り[intrusion=both]{ル}%
    \振り[intrusion=both]{多}{た}%
    \振り[intrusion=both]{少}{しょう}%
  }
\end{GCKEnv}

\bigskip

\begin{GCKEnv}{2\zw}
  \振り{所}{ゆ}\返り[intrusion=post]{二}\振り{\KanHyphen}{ゑ}\振り{以}{ん}

  \bigskip

  \parbox<t>{8\zw}{\GCKTateOn%
    \振り{雖}{いへど}\送り{モ}\\
    \振り{所}{ゆ}\返り[intrusion=post]{二}\振り{\KanHyphen}{ゑ}\振り{以}{ん}\\
    \振り{猶}{な}[ごと]\送り{ホ}[キヲ]
  \par}

  \bigskip

  \グ振り{所以}{ゆえん}\送り{ノミ}

  \bigskip

  \グ振り{所\返り[intrusion=post]{一}\KanHyphen 以}{ゆえん}\送り{ノミ}

  \bigskip

  % 生の \KanHyphen が直前の親文字へ重ならないことを確認する。
  \fbox{\グ振り{所\返り[intrusion=post]{一}\KanHyphen 以}{ゆえん}}

  \bigskip

  \グ振り[intrusion=both]{不可思議}{ふかしぎ}

  \bigskip

  \gckanbungroupruby{天地}{てんち}

  \bigskip

  \グ振り{不可思議}{ふしぎ}

  \bigskip

  \parbox<t>{6\zw}{\GCKTateOn%
    \グ振り{以心}{いしん}\グ振り{伝心}{でんしん}%
  \par}
\end{GCKEnv}

\end{document}
